%% expect: fallback fallback fallback fallback fallback fallback fallback fallback fallback native
%% reason 1: how the next token is read
%% reason 2: how the next token is read
%% reason 3: how the next token is read
%% reason 4: how the next token is read
%% reason 5: how the next token is read
%% reason 6: character constant
%% reason 8: assigned in a cell
\documentclass{article}
\usepackage{tikz-cd}
\begin{document}

% What comes before a token decides what it is, and the library has taken it for what it looks like.
\begin{tikzcd} A \noexpand\arrow[r] & B \end{tikzcd}
\begin{tikzcd} A \string\arrow[r] & B \end{tikzcd}
\begin{tikzcd} A \meaning\arrow[r] & B \end{tikzcd}
\begin{tikzcd} A \string& B \end{tikzcd}
\begin{tikzcd} A \expandafter\relax\arrow[r] & B \end{tikzcd}
\begin{tikzcd} \char`& A & B \end{tikzcd}
\begin{tikzcd} \char`\\ A & B \end{tikzcd}
\newcommand\q{}
\begin{tikzcd} A \let\q\arrow[r] & B \end{tikzcd}
\begin{tikzcd} A \afterassignment\arrow\count255=1 [r] & B \end{tikzcd}
% A plain \char constant, in the middle of a run.
\begin{tikzcd} \char`x A \arrow[r] & B \end{tikzcd}

\end{document}
